\documentclass[12pt]{article}

\usepackage[spanish]{babel}
\usepackage[utf8]{inputenc}
\usepackage[T1]{fontenc}
\usepackage{amsmath,amssymb}
\usepackage{geometry}
\usepackage{array}
\usepackage{enumitem}
\usepackage{lmodern}

\geometry{letterpaper, margin=2cm}

\begin{document}
	
	\begin{center}
		{\LARGE \textbf{Autoevaluación Evaluación II}}\\[0.2cm]
		{\large Contextualizada en Agronomía}
	\end{center}
	
	\vspace{0.4cm}
	
	\textbf{Instrucciones:} En cada caso realice todo el desarrollo en forma ordenada y no olvide responder la pregunta planteada.
	
	\vspace{0.5cm}
	
	\begin{enumerate}
		
		%-------------------------------------------------
		\item \textbf{De acuerdo con el último registro sobre los problemas más comunes observados en predios agrícolas, se tiene la siguiente información:}
		
		\begin{center}
			\begin{tabular}{|m{8cm}|>{\centering\arraybackslash}m{4cm}|}
				\hline
				\textbf{Problemas más comunes en el cultivo} & \textbf{Número de predios} \\
				\hline
				Daño por plagas & \\ \hline
				Enfermedades fúngicas & \\ \hline
				Déficit hídrico & \\ \hline
				Daño por heladas & \\ \hline
			\end{tabular}
		\end{center}
		
		\begin{enumerate}[label=\alph*)]
			\item Determine el número total de problemas más comunes observados en los predios agrícolas.
			\item Determine la diferencia entre los casos de daño por plagas y daño por heladas.
		\end{enumerate}
		
		%-------------------------------------------------
		\item \textbf{Resuelva la siguiente situación:} Para mejorar el rendimiento de un cultivo se deben fertilizar a
		\[
		(3n^2+5n)
		\]
		plantas y la dosis a aplicar a cada una es de
		\[
		(2n+7)\text{ mL}.
		\]
		Determine la cantidad total de mL que se deben emplear en la totalidad de la aplicación.
		
		%-------------------------------------------------
		\item \textbf{En un cultivo afectado por una plaga se prescribe un tratamiento fitosanitario.} Para el tratamiento completo se requieren \(x\) envases del producto, cada uno con
		\[
		(x^2-x+72)
		\]
		mL. Si se debe aplicar
		\[
		(x^2-9x)
		\]
		mL diarios, ¿cuántos días dura el tratamiento?
		
		%-------------------------------------------------
		\item \textbf{El administrador de un predio agrícola compró una cantidad de insumos de}
		\[
		Q=\frac{x^2-5x+6}{x-3}
		\qquad [\text{unidades}]
		\]
		a un valor de
		\[
		P=\frac{x^3+x^2+4x+4}{x^2-x-2}
		\qquad [\text{pesos por unidad}]
		\]
		El costo total \(C\) de los insumos es \(C=P\cdot Q\).
		
		\begin{enumerate}[label=\alph*)]
			\item Simplifique el costo total a su mínima expresión.
			\item Evalúe el costo total para \(x=1500\).
		\end{enumerate}
		
		%-------------------------------------------------
		\item \textbf{Un grupo de ingenieros agrónomos ha logrado medir la presencia en el suelo y en el ambiente de agentes que afectan el desarrollo de los cultivos. La fórmula encontrada es:}
		
		\[
		C=\left[
		\frac{x^2-x-12}{x^2-9}
		\cdot
		\frac{x^3+x^2-2x}{x^2-2x-8}
		\div
		\frac{x-1}{2x^2-6x}
		\right]
		\]
		
		Aplicando factorización y simplificación, encuentre la expresión más reducida que representa la fórmula para \(C\).
		
		\vspace{0.3cm}
		
		Simplifique además la siguiente expresión:
		
		\[
		\frac{\bigl[(8x^2-5)+4\bigr]^2(2x+7)^2-(4x-7)(3x-2)(3x+4)}
		{(5x+3)(3+5x)-9}
		\]
		
	\end{enumerate}
	
\end{document}